\section{An explicit differentiated clock}\label{sec:explicit}
The continuum mean $Lk_h/\pi$ suppresses the lattice phase. Its static
accuracy does not by itself control a time derivative on a growing
height window. We retain that phase exactly and replace the source mean
and its drift by explicit saddle quantities, with a separate error for
each replacement. Throughout this section $\ell,h,Y$ are fixed,
$Y>h$, and there is no time average.

Use the large saddle in \eqref{eq:saddle} and set
\begin{equation}
 M_0=2\pi^2e^{q_t(L)}\sqrt{2\pi/\kappa},\qquad
 \epsilon_+(t,x)=M_0\kappa^{-1}e^{-x^2/(2\kappa)},\qquad
 p=\pi/L.
 \label{eq:positiveRegulator}
\end{equation}
Implicit differentiation gives $L'=2L/\kappa$ and
$\kappa'=8L+16tL/\kappa-2$.
The candidate below uses only $L,\kappa$, the kernel widths and
fractional parts. Complete-source moments enter its proof and error
allowances, not its evaluated formula.

\subsection{A score estimate for the mean and drift}
\begin{lemma}[Complete-source score refinement]\label{lem:score}
Under the source gates of Appendix~\ref{app:source}, write
$\alpha=\mu'/\mu$ and put
\begin{align}
 M_-&=2\pi^2e^{q_t(L)}Z_-/\sqrt\kappa,&
 \overline p_0&=2\pi^2C_be^{-\pi}/M_-,& a&=1+2t/\kappa,\notag\\
 e_s&=2aA_2(4)+(5+\overline p_0)/\kappa,&
 e&=\min\{e_\mu,e_s\},\notag\\
 E_v^s&=2aA_3(4)+5A_1(0)/\kappa
                   +L\overline p_0/\kappa+e^2,\notag\\
 B_3^s&=4\{A_3(0)+e^3A_0(0)\},&
 E_\alpha&=2E_v^s+B_3^s/(L-e).\label{eq:scoreallowances}
\end{align}
If $L>e$, then
\begin{equation}
 |\mu-L|\le e=O(\kappa^{-1}),\qquad
 |v-\kappa^{-1}|\le E_v^s=O(\kappa^{-3/2}),\qquad
 |\alpha-2/\kappa|\le E_\alpha=O(\kappa^{-3/2}).
 \label{eq:scorebounds}
\end{equation}
The weight $4$ in $A_j(4)$ is a source weight, independent of the
physical transverse strip.
\end{lemma}
\begin{proof}
The complete prefactor in \eqref{eq:q} satisfies
$b\ge1-3/(2\pi)>1/2$ and $|b'/b|\le5$ on $u\ge0$.
Indeed the first term contributes at most $6/\pi$ to $|b'|$;
the remaining terms are bounded by
\[
 \sum_{n\ge2}(2n^2+16n^6)e^{-3(n^2-1)}<0.13.
\]
For this tail, its first term is $1032e^{-9}$ and the subsequent
ratios are less than $10^{-5}$. This also justifies differentiation
of the full series. Thus $|b'|<2.13$ and the stated score bound follows.

Let $\zeta=u-L$, $d=\E\zeta$, $\beta=b'/b$ and
$p_0=\Phi(0)/M$. The logarithmic source derivative is
\begin{equation}
 \partial_u\log(dP_t/du)=-\kappa\zeta+E_q(\zeta)+\beta(u),
 \qquad E_q=-\frac{\kappa+2t}{4}(e^{4\zeta}-1-4\zeta).
 \label{eq:scoreidentity}
\end{equation}
The exponential remainder gives
$|E_q|\le2(\kappa+2t)\zeta^2e^{4|\zeta|}$.
Integration by parts on the half-line, including its endpoint, yields
\begin{equation}
 \kappa d=\E E_q+\E\beta+p_0,\qquad
 \kappa\E\zeta^2=1-Lp_0+\E(\zeta E_q)+\E(\zeta\beta).
 \label{eq:scoreparts}
\end{equation}
Here $\int P_t'=-p_0$ and
$\int(u-L)P_t'=Lp_0-1$ fix the endpoint signs.
Apply the weighted moment bounds and $p_0\le\overline p_0$,
then use $v=\E\zeta^2-d^2$.
Finally differentiation of the source law gives
$\mu'=\Cov(u,u^2)=2\mu v+m_3$, where
$m_j=\E(u-\mu)^j$, and $|m_3|\le B_3^s$.
This proves the finite bounds. Their orders follow from
$A_j(4)=O(\kappa^{-j/2})$ and the superalgebraically small endpoint term.
No differentiated asymptotic remainder is used.
\end{proof}

\subsection{The lattice phase as a finite polynomial formula}
Write $\mathcal B_m(v)=B_m(\{v\})$ for the periodic Bernoulli
polynomial, with convention $B_1(v)=v-1/2$, and put
\[
 \theta_-=(x-\ell)/p+\tfrac12,\qquad
 \theta_+=(x+\ell)/p+\tfrac12.
\]
Use the interior endpoint derivatives of the polynomial kernel:
\begin{equation}
 a_j=\frac{315\sigma_j}{256\ell^{j+1}},\qquad
 (\sigma_4,\sigma_5,\sigma_6,\sigma_7,\sigma_8)
       =(384,3840,17280,40320,40320),\qquad
 b_j=\ell a_j+j a_{j-1},
 \label{eq:phasecoefficients}
\end{equation}
with other $a_j$ zero. Define
\begin{align}
 P(p,x)&=\frac1p+\sum_{j=4}^8\frac{p^ja_j}{(j+1)!}
 [\mathcal B_{j+1}(\theta_-)-(-1)^j\mathcal B_{j+1}(\theta_+)],
 \label{eq:finiteP}\\
 D(p,x)&=\sum_{j=4}^9\frac{p^jb_j}{(j+1)!}
 [\mathcal B_{j+1}(\theta_-)-(-1)^{j+1}\mathcal B_{j+1}(\theta_+)].
 \label{eq:finiteD}
\end{align}
The symbol $D(p,x)$ in these formulas is a lattice first moment;
it is distinct from the field denominator in \eqref{eq:definitions}.

\begin{proposition}[Exact phase and continuity]\label{prop:phase}
For $d_j=(j+1/2)p$,
\begin{equation}
 P=\sum_{j\in\mathbb Z}k_\ell(x-d_j),\qquad
 D=\sum_{j\in\mathbb Z}(x-d_j)k_\ell(x-d_j).
 \label{eq:latticeidentity}
\end{equation}
Consequently
\begin{equation}
 R_L=k_h(y)P(p,x),\qquad
 J_L=-\frac2\kappa k_h(y)\{xP(p,x)-D(p,x)\}
 \label{eq:explicitcandidate}
\end{equation}
have $R_L\ge0$ and satisfy $\partial_tR_L+\partial_xJ_L=0$
exactly. The current vector is $(J_L,0)$.
\end{proposition}
\begin{proof}
For a polynomial $f$ supported on $[-\ell,\ell]$ whose endpoint
jets of orders zero through three vanish, integration by parts terminates:
\[
 \widehat f(\omega)=-\sum_{j=4}^{\deg f}
 \frac{f^{(j)}(\ell)e^{-i\omega\ell}
             -f^{(j)}(-\ell)e^{i\omega\ell}}{(i\omega)^{j+1}}.
\]
Its shifted periodization has Fourier coefficients
$p^{-1}(-1)^n\widehat f(2\pi n/p)$.
Insert the Bernoulli Fourier identity \cite{dlmf24} and take
$f(s)=k_\ell(s)$ and $f(s)=sk_\ell(s)$.
The zero coefficients are respectively $1/p$ and zero; parity gives
\eqref{eq:finiteP}--\eqref{eq:finiteD}.
The series used are absolutely convergent, including at edge arguments.
Thus the finite formulas are exact, with no phase truncation error.
The locally finite lattice sums are $C^3$ through kernel-edge crossings.
Since $p'=-2p/\kappa$ and
$\partial_x(xP-D)=-pP_p$, continuity follows, with no edge point masses.
\end{proof}

\subsection{Differentiated root motion and regulator transfer}
We record the analytical ingredients that connect the explicit formula
to the adaptive density; full finite constants are given in Note~16
\cite{n16}. All moments below refer to the complete source law.
The weighted bounds in Appendix~\ref{app:source} extend to order six.
For $j\ge1$ one may use
\[
 B_j(Y):=\E[|u-\mu|^je^{Y|u-\mu|}]
 \le2^{j-1}e^{Ye_\mu}\{A_j(Y)+e_\mu^jA_0(Y)\}.
\]
The Gaussian fourth-moment comparison and centering give
$|m_4-3v^2|\le E_4=O(\kappa^{-9/4})$, with an explicit allowance
from the same $D_0,D_2,D_4$ bounds (Note~16, AT3--AT4).
Differentiating the source law, rather than a moment allowance, gives
\begin{equation}
 \mu'=2\mu v+m_3,\quad v'=2\mu m_3+m_4-v^2,\quad
 m_3'=2\mu(m_4-3v^2)+m_5-4vm_3.
 \label{eq:momentderivatives}
\end{equation}
In particular $|m_3'|=O(L\kappa^{-9/4})$.

Use \eqref{eq:carriergates} and \eqref{eq:growing}, and write the
root-location radius as $r_0=s/\mu$. For $f=H/M$ on the real axis,
Taylor expansion of the differentiated complete transform gives
\begin{align}
 f_z&=-\mu A\sin(\mu z)-vz\cos(\mu z)+\mathcal R_z,
 &|\mathcal R_z|&\le E_1,\notag\\
 E_1&=\tfrac12X^2B_3(0)+\tfrac16X^3\{\mu B_3(0)+B_4(0)\},
 \label{eq:derivativegate}\\
 d_z&:=\mu a_*\cos s-vX\sin s-E_1>0.\notag
\end{align}
The last strict inequality holds eventually and bounds $|f_z|$ at
each covered zero from below. Separately, cubic expansion in the
source covariance $f_t=\Cov(u^2,\cos(zu))$ yields
\begin{equation}
 f_t-\alpha zf_z=\tfrac16m_3'z^3\sin(\mu z)
 +(\alpha v-v'/2)z^2\cos(\mu z)
 +\mathcal R_t-\alpha z\mathcal R_z,
 \label{eq:driftcancellation}
\end{equation}
where
$|\mathcal R_t|\le T_4=X^4\{2\mu B_5(0)+B_6(0)+vB_4(0)\}/24$.
The cubic coefficient is $m_3'$, since
$2\mu m_4+m_5-vm_3-3\mu'v=m_3'$.
At a root $|\cos(\mu a_j)|\le\sin s$; hence
\begin{equation}
 |a_j'+\alpha a_j|\le E_{\rm vel}:=
 \frac{|m_3'|X^3/6+(|\alpha|v+|v'|/2)X^2\sin s
                       +T_4+|\alpha|XE_1}{d_z}
 =O(\kappa^{-7/6}).
 \label{eq:velocitybound}
\end{equation}
The moment identities and upper allowances replace the absolute
moment derivatives here by finite source bounds whenever desired.
This estimate holds uniformly over the growing zero population.

For regularization use wider collars of radius $R_c=1/(4\mu)$,
retaining $r_0$ for location accuracy, and strengthen the gate to
$s<1/8$. Put $\beta=1/4$ and
\[
 b=a_*\{e^{-\beta}\sin\beta-\delta\}>0,\qquad c=b/(2R_c).
\]
Then $|H/M|\ge be^{\mu|y|}$ off the collars and on their boundaries,
and $|H/M|\ge c|z-a_j|$ inside them.
There are at most $n=2\mu(\ell+R_c)/\pi+2$ relevant collars.
For $\eta=\sqrt{2\pi}/(Z_-\kappa)$ and $a_e=\eta/c$, set
\begin{align}
 W_0&=4\ell h\,\eta^2/b^2
              +n\pi a_e^2\log(1+4R_c^2/a_e^2),&
 \Lambda_0&=W_0+n\pi a_e^2,\label{eq:collarintegrals}\\
 Q_2&=\mu^2B_0(Y)+2\mu B_1(Y)+B_2(Y),&
 V&=|\alpha|X+E_{\rm vel},\notag\\
 I_0&=\left\{Q_2/b+\frac{V}{R_c-r_0}\right\}W_0
                                      +n\pi^2Va_e.\label{eq:poleintegral}
\end{align}
Here $B_0(Y)=\E e^{Y|u-\mu|}$; the moment $B_2(Y)$ differs from
the global source derivative constant $B_2$ in Lemma~\ref{lem:B2}.
On a full kernel support $\Omega$,
\begin{equation}
 \int_\Omega w\le W_0=O(\kappa^{-2}),\quad
 \int_\Omega\Lambda_q\le\Lambda_0=O(\kappa^{-2}),\quad
 \int_\Omega|H_{zz}/H|w\le I_0
             =O(L^2\kappa^{-2}+\kappa^{-5/3}).
 \label{eq:weakorders}
\end{equation}
For the first two bounds integrate the rational and logarithmic
simple-zero profiles over disks of radius $2R_c$ about the actual zeros.
For the third, $H_{zz}/H$ has residue $a_j'$ at a covered zero.
Subtract that pole and use the boundary maximum principle for its
regular part. The pole contributes at most
$V\int w/|z-a_j|\le\pi^2Va_e$ per collar.
This step preserves the small pole residue rather than bounding the
meromorphic quotient uniformly across a zero.

Let $J_\nu=\sum_j a_j'\delta_{a_j}$. With
$\partial_{\bar z}=(\partial_x+i\partial_y)/2$, one has distributionally
\begin{equation}
 j-J_\nu=\frac1\pi\left\{\frac{H_{zz}}H w(1-w)\partial_{\bar z}r
             -\partial_{\bar z}\left(\frac{H_{zz}}H w\right)\right\}.
 \label{eq:weakcurrentbridge}
\end{equation}
This follows by differentiating $w$ off the zeros and using
$\partial_{\bar z}(1/(z-a))=\pi\delta_a$ at each simple zero.
For norms on the full buffer set
\begin{align*}
 b_1&=\|\nabla r\|_\infty,&b_2&=\|\Delta r\|_\infty,&a_0&=\|r_t\|_\infty,\\
 b_{1t}&=\|\nabla r_t\|_\infty,&b_{2t}&=\|\Delta r_t\|_\infty,\\
 \mathcal K_r(\psi)&=\|\Delta\psi\|_\infty+2b_1\|\nabla\psi\|_\infty
                     +(b_2+b_1^2)\|\psi\|_\infty,\\
 \mathcal K_{r,t}(\psi)&=2b_{1t}\|\nabla\psi\|_\infty
                     +(b_{2t}+2b_1b_{1t})\|\psi\|_\infty.
\end{align*}
Equations \eqref{eq:adaptiveWeak} and \eqref{eq:weakcurrentbridge}
then give the finite allowances
\begin{align}
 E_R(\psi)&=\mathcal K_r(\psi)\Lambda_0/(4\pi),\notag\\
 E_J(\psi)&=(\|\nabla\psi\|_\infty+b_1\|\psi\|_\infty)I_0/(2\pi),\notag\\
 E_T(\psi)&=\{\mathcal K_r(\psi)(a_0W_0+2I_0)
                          +\mathcal K_{r,t}(\psi)W_0\}/(4\pi).
 \label{eq:transportallowances}
\end{align}
They bound, respectively, the pairings of $\rho-\nu$,
$j-J_\nu$, and $\partial_t\rho-\partial_t\nu$ against $\psi$.
For the last estimate differentiate the exact weak identity, using
\[
 (\Lambda_q)_t=w(r_t+2\Re(H_{zz}/H)),\qquad
 w_t=w(1-w)(r_t+2\Re(H_{zz}/H)).
\]
The simple-pole singularities are locally integrable; local simple-zero
factorization justifies the distributional time derivative.
For \eqref{eq:positiveRegulator},
\begin{gather*}
 a_0\le2L^2+3|\kappa'|/\kappa+X^2|\kappa'|/\kappa^2,\qquad
 b_1\le2X/\kappa,\quad b_2=2/\kappa,\\
 b_{1t}\le2X|\kappa'|/\kappa^2,\qquad b_{2t}=2|\kappa'|/\kappa^2.
\end{gather*}
Thus $E_R(K)=O(\kappa^{-2})$ and
$E_J(K),E_T(K)=O(\kappa^{-5/3}+L^2\kappa^{-2})$.
The signed weak source $K*\mathcal S_\epsilon$ has the latter order
by the two continuity equations. Also
$|K*\mathcal B_\epsilon|=O(L^2\kappa^{-2})$.
These are bounds on the actual correlated terms; they are not products
of separately averaged factors.

\subsection{The explicit approximation theorem}
Let $k_i=\|k_\ell^{(i)}\|_\infty$ and $h_0=\|k_h\|_\infty$.
Matching zeros to the mean clock $c_j=(j+1/2)\pi/\mu$ first gives
finite current and time-derivative errors
\begin{align}
 A_J&=E_J(K)+nh_0\{E_{\rm vel}k_0+|\alpha|r_0(k_0+Xk_1)\},\notag\\
 A_T&=E_T(K)+nh_0\{E_{\rm vel}k_1+|\alpha|r_0(k_1+Xk_2)\}.
 \label{eq:meanclockcosts}
\end{align}
Both are $O(L\kappa^{-7/6})$ by the preceding estimates.

\begin{theorem}[Explicit density and transport]\label{thm:explicit}
Take \eqref{eq:positiveRegulator}, \eqref{eq:kernel} and
\eqref{eq:growing}. Impose the source, carrier and derivative gates
above, $s<1/8$, and the full buffers
\begin{equation}
 |x|+\ell+2R_c+\delta_c<X,\qquad |y|+h+2R_c<Y,\qquad
 \delta_c=Xe/L.
 \label{eq:explicitbuffer}
\end{equation}
Put $X_*=X+\delta_c$ and
$n_*=2\mu(\ell+R_c+\delta_c)/\pi+2$.
Then the density error is at most
\begin{equation}
 \mathcal E_R=E_R(K)+h_0n_*k_1(r_0+\delta_c),
 \label{eq:finalRerror}
\end{equation}
and the current and density time-derivative errors are at most
\begin{align}
 \mathcal E_J&=A_J+h_0n_*
       \{E_\alpha Xk_0+(2/\kappa)\delta_c(k_0+X_*k_1)\},\notag\\
 \mathcal E_T&=A_T+h_0n_*
       \{E_\alpha Xk_1+(2/\kappa)\delta_c(k_1+X_*k_2)\}.
 \label{eq:finalTransportError}
\end{align}
In particular, uniformly on these growing windows as $t\to\infty$,
\begin{align}
 \|K*\rho-R_L\|_\infty&=O(\kappa^{-1/2}),\label{eq:explicitDensityRate}\\
 \|K*\mathbf j-(J_L,0)\|_\infty+
 \|\partial_t(K*\rho)-\partial_tR_L\|_\infty
 &=O(L\kappa^{-7/6})=o(L/\kappa).
 \label{eq:explicitTransportRate}
\end{align}
All gates hold eventually; no finite certified onset is asserted.
\end{theorem}
\begin{proof}
For every relevant matched label,
$|c_j-d_j|=|c_j||\mu-L|/L\le\delta_c$.
Matching $a_j$ to $c_j$ and then to $d_j$ costs
$k_1(r_0+\delta_c)$ per density term. At most $n_*$ labels enter;
the buffers include their full collars.
For current compare $\alpha c_jk_\ell(x-c_j)$ with
$(2/\kappa)d_jk_\ell(x-d_j)$.
The drift error costs $E_\alpha Xk_0$, and the position error costs
$(2/\kappa)\delta_c(k_0+X_*k_1)$.
The same argument with $k_\ell'$ gives the time-derivative comparison.
The mean-clock errors \eqref{eq:meanclockcosts} follow from
\eqref{eq:velocitybound} and the weak allowances.
Zero extension of the $C^3$ kernel validates these comparisons at
kernel-edge crossings.

Now $n_*=O(L)$, $r_0=O(\kappa^{-1/2}/L)$ and
$\delta_c=O(\kappa^{-2/3}/L)$.
The new drift error costs $O(L\kappa^{-7/6})$, and the position error
in transport is $O(\kappa^{-4/3})$.
Combine these with \eqref{eq:transportallowances} and
\eqref{eq:meanclockcosts}. The finite phase formulas add zero error.
All appearances of $\mu$ and positive moments in the allowances can be
bounded using the source majorants, with $L-e$ in denominators and
$L+e_\mu$ in upper counts.
\end{proof}
This theorem controls the unnormalized density and its transport.
It does not assign a potential or the actual theta covariance to
$(R_L,J_L)$ by forcing an equation to close. The distinction matters
because the true adaptive continuity law has the small source in
\eqref{eq:adaptiveCurrent}, whereas the explicit reference is exactly
conservative.
